\section{Finite-frequency amplitude and phase certificates}
The positive transport action now has an explicit finite-frequency estimate.
Write $B=b_{\rm cos}\sqrt Q>0$, $U=B''/B\in L^1[a,b]$, $c_T=1$, and retain the same incoming initial data for the canonical mode and its plane-wave comparator:
\[
 v(a)=1,\quad v'(a)=-\ii k,\qquad
 v_0(\eta_c)=e^{-\ii k(\eta_c-a)},\quad k>0.
\]
Variation of constants gives the exact equation
\[
 v=v_0+\int_a^{\eta_c}\frac{\sin[k(\eta_c-s)]}{k}U(s)v(s)\,\dd s.
\]
Set $M(\eta_c)=k^{-1}\int_a^{\eta_c}|U|\,\dd s$ and $\rho=e^M-1$. Gronwall, followed by substitution and differentiation, gives
\[
 |v-v_0|\le\rho,\qquad |v'-v_0'|\le k\rho.
\]
Indeed $|v(s)|\le e^{M(s)}$ and the resulting integral is exactly $e^{M(\eta_c)}-1$. No derivative bound on $U$ is needed. For $h=v/B$ and $h_0=v_0/B$,
\[
\begin{split}
 |h-h_0|&\le\rho/B,\\
 |h'-h_0'|&\le(k+|B'/B|)\rho/B.
\end{split}
\]
If $\rho<1$, the same disk gives
\[
 \big||v|-1\big|\le\rho,\qquad
 |\arg(v/v_0)|\le\arcsin\rho,
\]
with one continuous phase unwrapping from $a$. The real potential preserves $\operatorname{Im}(\bar v v')=-k$, so the actual phase rate is $-k/|v|^2$ and
\[
 \left|\frac{\dd\arg v}{\dd\eta_c}+k\right|
 \le k\big[(1-\rho)^{-2}-1\big].
\]
The positive factor $B$ changes no phase. These are separate amplitude, accumulated-phase and phase-rate certificates.

For a fully specified declared example, take $0\le\eta_c\le1$ and
\[
 b_{\rm cos}=\frac{10+\eta_c}{11},\qquad
 Q=\left(\frac{10+\eta_c}{10}\right)^2,\qquad k=100.
\]
Then $B=(10+\eta_c)^2/110$, $U=2/(10+\eta_c)^2$ and
\[
 M(1)=\frac1{5500},\qquad \rho(1)<\frac1{5499}.
\]
The last rational bound uses $e^x\le(1-x)^{-1}$ for $0\le x<1$. The leading endpoint factor is $B_E/B_R=100/121$, and the propagation-only distance ratio is $\sqrt{Q_R/Q_E}=11/10$. The finite-frequency amplitude disk is below $1.819\times10^{-4}$. This receiving example assigns no microscopic SAM $Q$ or physical conformal-time scale. An elementary exact basis,
\[
 w(\eta_c)=e^{-\ii k\eta_c}
 \left(1-\frac{\ii}{k(10+\eta_c)}\right),\quad\bar w,
\]
with Wronskian $2\ii k$, supplies the independent matched-initial-data check.

\section{Generation, propagation and detector normalization}
A source term $\mathcal J h$ in one normalized Fourier mode gives
\[
 (B^2h')'+B^2k^2h=\mathcal J,\qquad
 v''+(k^2-U)v=\mathcal J/B.
\]
For a zero incoming incremental mode and impulse $\mathcal J=j\delta(\eta_c-E)$, the canonical momentum jumps by $j$. If $G(R,E)$ has unit initial velocity and zero initial displacement, then
\[
 h(R)=\frac{jG(R,E)}{B_EB_R},\qquad
 W_{\rm impulse}=\frac{j^2}{2B_E^2}.
\]
Thus an equal source impulse and an equal emitted field preparation are different protocols when $Q_E$ differs.

Keep the preceding $b_{\rm cos}$ but set $Q=4[(10+\eta_c)/10]^2$. This leaves $U$ unchanged; now $Q_E=4$ and $Q_R=121/25$. Compare an electromagnetic mode with $Q=1$, the same $k=100$, and $E=0$, $R=1$. For $x_E=10$, $x_R=11$ and $L=1$, the complete retarded Green functions are
\[
\begin{split}
 G={}&\left(1+\frac1{k^2x_Ex_R}\right)\frac{\sin kL}{k}\\
 &+\left(\frac1{x_R}-\frac1{x_E}\right)\frac{\cos kL}{k^2},\\
 G_{\rm EM}={}&\frac{\sin kL}{k}.
\end{split}
\]
Here $G/G_{\rm EM}=1.0001557233563114\ldots$ at this specified nonzero readout. The absolute Green error also satisfies $k|G-G_{\rm EM}|\le e^M-1$; no ratio is assigned at a zero of $G_{\rm EM}$.

Let $d$ be the receiving detector-gain ratio. The recorded amplitudes obey exactly
\[
 \frac{A_T}{A_{\rm EM}}=
 d\frac{j_T}{j_{\rm EM}}\frac1{\sqrt{Q_EQ_R}}
 \frac{G}{G_{\rm EM}}.
\]
Equal impulses and $d=1$ give $(5/22)(G/G_{\rm EM})$. Matching the emitted $h'$ requires $j_T=Q_Ej_{\rm EM}=4j_{\rm EM}$ and gives $(10/11)(G/G_{\rm EM})$. A detector gain $d=3/2$ changes the latter to $(15/11)(G/G_{\rm EM})$. The inferred amplitude distances are the reciprocal ratios. Only the matched-emission, matched-detector protocol has the leading propagation-only ratio $11/10$.

For $j_{\rm EM}=1$, the electromagnetic impulse costs $121/200$ per normalized mode; the equal-impulse tensor mode costs $121/800$; matched emission costs $121/50$. During propagation the prescribed background also does signed work:
\[
 \frac{\dd H}{\dd\eta_c}
 =BB'k^2h^2-\frac{B'}{B^3}\Pi^2
\]
for each real mode outside the impulse. Source preparation, prescribed-background work and a detector's absorbing or disturbing measurement remain in their respective accounts. The displayed linear detector gain is a declared receiving convention.

\section{The original causal tensor carrier as a signal map}
The next calculation returns to the actual retained quadrupole TT carrier, distinct from the preceding electromagnetic and cosmological adapters. Keep its two signed source quadratures, $\kappa=4/5$, $c_*=1$, regular interior, full outgoing exterior and original radius $R$. For $r=\varrho R$, $\varrho\ge1$, let
\[
\begin{split}
 s(w)&=(r+Rw)/c_*,\\
 K_{\varrho}(w)&=\frac1{\kappa\varrho}
 P_2\!\left(-w+\frac{1-w^2}{2\varrho}\right).
\end{split}
\]
The complete driven receiver is
\[
 x_{\rm out}(t)=\int_{-1}^{1}K_{\varrho}(w)f(t-s(w))\,\dd w.
\]
A homogeneous incoming/prepared field adds its own receiver term. The retained maintained-pulse experiment has zero incremental incoming field on the same stationary reference; the full background and field energy remain. The frequency response is
\[
 \chi(\omega)=\frac{2\ii z}{\kappa}j_2(z)h_2^{(1)}(z\varrho),
 \qquad z=\omega R/c_*.
\]
All near, induction and radiation terms are retained. The first and last kernel supports are $(r-R)/c_*$ and $(r+R)/c_*$. On the original small-frequency branch $j_2(z)>0$, the group delay is
\[
 \tau=\frac r{c_*}
 \frac{u^4}{u^4+3u^2+9},\qquad u=\omega r/c_*.
\]
A signed kernel can have a group delay outside its causal support interval; the two observables remain distinct.

For a carrier-modulated envelope $f=a(t)e^{-\ii\omega_0t}$, choose $\chi_0=\chi(\omega_0)\ne0$ and $\tau=(\dd\arg\chi/\dd\omega)|_{\omega_0}$. The gain/delay approximation is
\[
 x_{\rm app}=\chi_0e^{-\ii\omega_0t}a(t-\tau).
\]
More explicitly, write $\tau=(\dd/\dd\omega)\arg\chi|_{\omega_0}$ and define
\[
\begin{split}
 L_1&=|\chi'(\omega_0)-\ii\tau\chi_0|,\\
 D_2&=\int |K_{\varrho}(w)|[s(w)-\tau]^2\,\dd w.
\end{split}
\]
Taylor's remainder for $e^{\ii\nu(s-\tau)}$ gives
\[
 |\chi(\omega_0+\nu)-\chi_0e^{\ii\nu\tau}|
 \le |\nu|L_1+\tfrac12\nu^2D_2.
\]
If the envelope spectrum lies in $|\nu|\le\beta$, Plancherel gives
\[
 \|x_{\rm out}-x_{\rm app}\|_2
 \le(\beta L_1+\tfrac12\beta^2D_2)\|a\|_2.
\]
The first term is amplitude dispersion: at the stated group delay $L_1=|\dd|\chi|/\dd\omega|$. For a bounded smooth envelope, direct Taylor expansion instead gives the pointwise bound
\[
 |x_{\rm out}-x_{\rm app}|
 \le L_1\|a'\|_\infty+\tfrac12D_2\|a''\|_\infty.
\]
This version needs no frequency truncation, and applies to a compact pulse.

\section{A complete source pulse and a nonzero phase event}
Retain the actual four-cycle source
\[
 f_{ab}(t)=g_{ab}A(\Delta t)e^{-\ii\Delta t},\qquad
 \Delta=\frac2{573049},
\]
where $A(\theta)=\sin^4(\theta/8)$ on $[0,8\pi]$ and zero outside. The six signed $g_{ab}$, equal pair weights $1/6$, both tensor polarizations, four undriven source diagonal states and complete maintenance-work account are unchanged. The envelope is $C^3$ with bounded fourth weak derivative.

The complete signed delay moments satisfy
\[
\begin{split}
 m_0&=\frac2{5\kappa\varrho^3},\quad m_1=m_3=0,\\
 \frac{m_2}{m_0}&=\frac{R^2}{7c_*^2}-\frac{r^2}{3c_*^2},\quad
 m_n=\int K_{\varrho}s^n\,\dd w.
\end{split}
\]
These follow by expanding the full argument of $P_2$ and integrating monomials. At $\varrho=2$, with $(R/c_*)^n$ factored out, moments zero through four are
\[
 (1/16,\ 0,\ -25/336,\ 0,\ 313/336).
\]
Thus odd-moment cancellation and a negative signed second moment quantify the near-field behavior.

Taylor's formula, including both compact-pulse edges, gives
\[
\begin{split}
 |x_{\rm out}-m_0f|
 &\le\frac{|m_2|}2\|f''\|_\infty
       +\frac{D_4}{24}\|f^{(4)}\|_\infty,\\
 D_4&=\int|K_{\varrho}|s^4\,\dd w
 \le\frac{2(r+R)^4}{\kappa\varrho c_*^4}.
\end{split}
\]
For $\varrho\ge1$, the Legendre argument maps $[-1,1]$ monotonically onto $[1,-1]$, so $|P_2|\le1$. The exact five-frequency interior expansion gives
\[
 \|f''\|_\infty\le|g|\Delta^2\frac{17}{16},\qquad
 \|f^{(4)}\|_\infty\le|g|\Delta^4\frac{709}{512}.
\]
The zero extension has the same weak-derivative bounds. No frequency or radial tail is cut.

At the original receiver $r=R+64$, use the complete authenticated Chapter I radius ball, rather than treating a printed decimal as exact. Outward rational arithmetic yields, for every time and every original signed lane,
\[
 \frac{|x_{\rm out}-m_0f|}{m_0|g|}<8.994\times10^{-9}.
\]
Also $|\chi(\Delta)-m_0|/m_0<8.440\times10^{-9}$ and
\[
 0<\tau<1.792884\times10^{-14}
\]
in original model-time units, while first causal arrival remains exactly $64$. Replacing the static gain by the full monochromatic gain and its group delay gives
\[
 \frac{\left|x_{\rm out}-g\chi(\Delta)e^{-\ii\Delta t}
 A[\Delta(t-\tau)]\right|}{m_0|g|}
 <1.744\times10^{-8}.
\]
This is a uniform absolute-envelope error normalized by peak static gain. A phase disk requires a positive received envelope: if $A[\Delta(t-\tau)]\ge a_*>0$, its radius is at most
\[
 \arcsin\!\left[\frac{\epsilon_{\rm pulse}}
 {a_*(1-\epsilon_{\rm gain})}\right],
\]
when the argument is below one. It supplies no phase statement at zeros.

A sharper phase certificate is available at the actual peak event $t_*=4\pi/\Delta$. Every retarded source argument then lies inside the smooth pulse. After demodulation and division by signed $g$,
\[
 Y=\int K_{\varrho}(w)e^{\ii\Delta s}
       \cos^4(\Delta s/8)\,\dd w.
\]
The sine expansion has positive weights and relative frequencies
\[
 \begin{aligned}
 \text{Weights}:\;&(3/8,1/4,1/4,1/16,1/16),\\
 \text{Frequency ratios}:\;&(1,5/4,3/4,3/2,1/2).
 \end{aligned}
\]
Their fifth-derivative sum is $857/512$. Since $m_1=m_3=0$,
\[
 \begin{aligned}
 |\operatorname{Im}Y|&\le\frac{857}{512}\frac{\Delta^5D_5}{120},\\
 |\operatorname{Im}\chi|&\le\frac{\Delta^5D_5}{120},
 \end{aligned}
\]
where $D_5\le2(r+R)^5/(\kappa\varrho c_*^5)$. The preceding real-part lower bounds then give
\[
\begin{split}
 |\arg Y-\arg\chi|\le{}&\frac{\Delta^5D_5}{120m_0}
 \left[\frac{857/512}{1-\epsilon_{\rm static}}
       +\frac1{1-\epsilon_{\rm gain}}\right]\\
 <{}&9.782\times10^{-15}\ \text{radians}.
\end{split}
\]
The gain/delay comparator's envelope is positive at this event and changes no phase. This is one nonzero-amplitude phase event; timing, incoming-field contamination, detector disturbance and statistical error remain in the measurement protocol. Source maintenance, free/drive/counterterm energy and complete emitted-field energy remain in the retained ledger.
